\section{IMPLEMENTATION DETAILS}
\label{sec:implementation_details}

\subsection{Data Preprocessing}
\label{sec:data_preprocessing}

We preprocess the BraTS data at the patient level, where each subject contains multi-modal 3D NIfTI volumes. In our current local copy, the dataset contains 625 subjects. For each subject, we load the T2-weighted volume (\texttt{t2w}) as the reconstruction target modality and the native T1-weighted volume (\texttt{t1n}) as the auxiliary anatomical modality. Both volumes are converted to \texttt{float32} and transposed into a slice-first representation of shape $(D, H, W)$.

Each 3D volume is independently normalized to $[0,1]$ using min-max normalization. We further discard nearly empty boundary slices using a low-variance threshold.

To emulate accelerated MRI acquisition, we transform each fully sampled T2 slice to k-space and apply a variable-density Cartesian undersampling mask. Unless otherwise stated, we use an acceleration factor of $R=5$ with center fraction $0.08$. The inverse FFT magnitude is used as the aliased T2 input.

For each retained slice, we store the aliased T2 image, the fully sampled T2 image, the fully sampled T1 image, the undersampling mask, and the undersampled complex k-space.

We split the dataset at the patient level to avoid leakage across slices from the same subject. Specifically, subjects are randomly shuffled with seed $42$ and divided into training, validation, and test sets with ratios of $70\%$, $15\%$, and $15\%$, respectively. Under the current dataset size of 625 subjects, this produces 437 training subjects, 93 validation subjects, and 95 test subjects.

\subsection{U-Net (Task2) Training Settings}
\label{sec:unet_training_settings}

For Task 2, we train a baseline single-modal U-Net reconstructor on the preprocessed dataset. The network takes the aliased T2 slice as input and predicts the fully sampled T2 slice in the image domain.

The baseline model is a standard encoder-decoder U-Net with depth $4$ and base width $16$.

We optimize the model using MSE loss and Adam with learning rate $1\times10^{-3}$. Training runs for $50$ epochs with batch size $128$, and validation is performed every epoch.

We use a short warmup stage followed by \texttt{ReduceLROnPlateau} learning-rate decay, together with gradient clipping and early stopping.

During evaluation, we report NMSE, PSNR, SSIM, and LPIPS.

We keep the best checkpoint according to validation loss.

\subsection{Multi-modal Unrolled Reconstruction Network (Task3) Training Settings}
\label{sec:multimodal_unrolled_training_settings}

For Task 3, we train a multi-modal unrolled reconstruction network on the same preprocessed dataset. The model takes the aliased T2 image, fully sampled T1 image, undersampling mask, and undersampled T2 k-space as inputs, and predicts the fully sampled T2 image.

The model uses a T1-guided unrolled architecture with $5$ iterations. Each iteration consists of a learned denoising block followed by a hard data-consistency step in k-space. In our main setting, the denoiser width is $32$ channels.

The training objective combines MSE loss, Haar high-frequency wavelet loss, and LPIPS perceptual loss, with weights $1.0$, $0.03$, and $0.1$, respectively. We train the model with Adam for $80$ epochs using batch size $32$ and gradient accumulation over $2$ steps, giving an effective batch size of $64$. The initial learning rate is $1.5\times10^{-4}$.

We use one-epoch warmup followed by \texttt{ReduceLROnPlateau} learning-rate decay, together with gradient clipping and mixed-precision training.

Validation is performed every epoch, and we report NMSE, PSNR, SSIM, and LPIPS.

We keep the best checkpoint according to validation loss.
